\documentclass[12pt,a4paper,twoside,openright]{book}

\usepackage[italian]{babel}
\usepackage[utf8]{inputenc}

\usepackage{style/isi_style_lt}

\usepackage{amsmath,amsfonts,amssymb,amsthm}
\usepackage{caption}
\usepackage[usenames]{color}
\usepackage{enumerate}
\usepackage{fancyhdr}
\usepackage{fancyvrb}
\usepackage{float}
\usepackage{graphicx}
\usepackage{indentfirst}
\usepackage{listings}
\usepackage{marvosym}
\usepackage{multicol}
\usepackage{sectsty}
\usepackage{subcaption}
\usepackage{tocloft}
\usepackage[table]{xcolor}
\usepackage{url}

\AtBeginDocument{%
	\renewcommand{\contentsname}{Indice}
	\renewcommand\tablename{Tabella}
	\renewcommand\figurename{Figura}
	\renewcommand{\lstlistingname}{Listato}
	\renewcommand{\refname}{Riferimenti}
}

\definecolor{dkgreen}{rgb}{0,0.6,0}
\definecolor{gray}{rgb}{0.5,0.5,0.5}
\definecolor{mauve}{rgb}{0.58,0,0.82}

\lstset{
  frame=single,
  captionpos=b,
  language=Java,
  aboveskip=3mm,
  belowskip=3mm,
  showstringspaces=false,
  columns=flexible,
  basicstyle={\small\ttfamily},
  numbers=none,
  numberstyle=\tiny\color{gray},
  keywordstyle=\color{blue},
  commentstyle=\color{dkgreen},
  stringstyle=\color{mauve},
  breaklines=true,
  breakatwhitespace=true,
  tabsize=3
}

\makeatletter
\def\cleardoublepage{
	\clearpage\if@twoside \ifodd\c@page\else
	\hbox{}
	\thispagestyle{empty}
	\newpage
	\if@twocolumn\hbox{}\newpage\fi\fi\fi
}

\makeatother

\setlength{\textwidth}{14cm}
\setlength{\textheight}{21cm}
\setlength{\footskip}{3cm}

\setlength{\hoffset}{0pt}
\setlength{\voffset}{0pt}

\setlength{\oddsidemargin}{1cm}
\setlength{\evensidemargin}{1cm}

\universita{Alma Mater Studiorum -- Università di Bologna}

\campus{Campus di Cesena}

\scuola{Scuola di Scienze}

\corsodilaurea{Corso di Laurea in Ingegneria e Scienze Informatiche}

\titolo{Qui il titolo\\della\\tesi}

\materia{Qui la materia}

\laureando{Gianni Cognome}

\relatore[Prof.]{Nome Cognome}
\correlatoreA[Ing.]{Nome Cognome}
%\correlatoreB[Ing.]{Nome Cognome}

\sessione{Seconda} % 'Prima', 'Seconda', 'Terza'

\annoaccademico{2018 -- 2019}

\parolechiave %devono essere 5 keywords
{prima key}
{seconda key}
{terza key}
{quarta key}
{quinta key}



\dedica{eventuale dedica} %commentare per non inserire la dedica

\makeindex

\begin{document}

\frontmatter 

\maketitle

\tableofcontents

\chapter{Introduzione}
\markboth{INTRODUZIONE}{INTRODUZIONE}

Qui il testo dell'introduzione alla tesi. Generalmente l'introduzione non dovrebbe superare le 2/3 pagine e dovrebbe essere scritta solo alla fine.
	
\mainmatter

\pagestyle{fancy} 
\fancyhead[LE,RO]{\thepage}
\fancyfoot{}

\chapter{Titolo del capitolo}
Qui una breve descrizione del contenuto del capitolo o, eventualmente, testo introduttivo alle sezioni che seguono

\section{Prima Sezione}
Qui il testo della prima sezione. Questa è  una parola in \textit{corsivo}, questa invece è una parola in \textbf{grassetto} e questa è una parola in \texttt{monotype}.

\begin{center}
Questo è un paragrafo centrato!
\end{center}

Dimensione del testo:\\

\LARGE{Testo} \Large{Testo} \large{Testo} \normalsize{Testo} \small{Testo} \footnotesize{Testo}\\

Qui una citazione bibliografica \cite{bib001}.

Qui un link \url{www.google.it}


\subsection{Prima Sottosezione}

Qui l'esempio di un elenco puntato:

\begin{itemize}
\item Primo elemento,
\item Secondo elemento,
\item Terzo elemento.
\end{itemize}

Qui l'esempio di un elenco con un sotto-elenco:

\begin{itemize}
\item Primo elemento,
\begin{itemize}
\item Primo elemento del sotto-elenco
\item Secondo elemento del sotto-elenco
\end{itemize}
\item Secondo elemento,
\item Terzo elemento.
\end{itemize}

Qui infine l'esempio di un elenco numerato

\begin{enumerate}
\item Testo
\item Testo
\item Testo
\end{enumerate}

\section{Seconda Sezione}\label{sec:Sezione2}

Qui l'esempio di una formula matematica:

$ G = \gamma\dfrac{m_{1}m_{2}}{r^{2}} $

\section{Terza Sezione}

Qui il riferimento ad una sezione utilizzando una label: Sezione \ref{sec:Sezione2}

Qui la definizione di due paragrafi con titoli.

\paragraph{Titolo 1} testo del paragrafo

\paragraph{Titolo 2} testo del paragrafo, con una nota a piè di pagina\footnote{Questo è il testo della nota, in cui si può \textit{utilizzare} qualunque stile di \textbf{formattazione}}


\chapter{Capitolo con Immagini e Codice}

Introduzione al capitolo

\section{Prima sezione del secondo capitolo}

Qui di seguito un'immagine. Tutte le immagini devono essere inserite nella cartella \texttt{images}

\begin{figure}[H]
\centering
\includegraphics[scale=0.8]{images/placeholder.jpg}
\caption{Testo della didascalia}
\label{fig:figure1}
\end{figure}

Qui invece il testo con il riferimento all'immagine: Figura \ref{fig:figure1}.

Qui un'immagine con due sotto-immagini:

\begin{figure}[H]
\centering
\begin{subfigure}[b]{0.42\textwidth}
\includegraphics[width=\textwidth]{images/placeholder.jpg}
\caption{Didascalia prima figura}
\label{fig:figure2}
\end{subfigure}
\qquad
\begin{subfigure}[b]{0.42\textwidth}
\includegraphics[width=\textwidth]{images/placeholder.jpg}
\caption{Didascalia seconda figura}
\label{fig:figure3}
\end{subfigure}
\caption{Didascalia globale}
\label{fig:figures12}
\end{figure}

Altri tipi di riferimenti: Figura \ref{fig:figures12} oppure \ref{fig:figure3}.

\section{Seconda sezione con codice sorgente}

Il codice sorgente può essere inserito mediante un link al file del codice sorgente, in questo caso il file con il codice va inserito nella cartella \texttt{code}.

\lstinputlisting[caption={Didascalia o il nome del file}, label=lst:listato1]{code/Main.java}



\chapter{Titolo}

\section{Titolo della sezione}

\subsection{Titolo della sottosezione}

\chapter*{Conclusioni}
\addcontentsline{toc}{chapter}{Conclusioni}
\markboth{CONCLUSIONI}{CONCLUSIONI}

Qui il testo delle conclusioni alla tesi. Non deve essere un riepilogo di quanto fatto nella tesi ma piuttosto le conclusioni raggiunte relative al lavoro svolto.
	
\chapter*{Ringraziamenti}
\addcontentsline{toc}{chapter}{Ringraziamenti}
\markboth{RINGRAZIAMENTI}{RINGRAZIAMENTI}

Qualora lo si desideri è possibile inserire qui il testo dei ringraziamenti alle persone che hanno contribuito in qualche modo al percorso che ha portato al lavoro di tesi.
	
\backmatter	

\begin{thebibliography}{99}
\addcontentsline{toc}{chapter}{Bibliografia}

\bibitem{bib001} Autore1, Autore2,
{\em Titolo dell'opera},
Editore, Anno.

\bibitem{bib002} Autore1, Autore2,
{\em Titolo dell'opera},
Editore, Anno.

\bibitem{bib003} Autore1, Autore2,
{\em Titolo dell'opera},
Editore, Anno.

\end{thebibliography}

\end{document}